% ###################    CONFIGURACIÓN ################### 

% Configuración de estilos generales de la página
\documentclass[12pt,
               % twocolumn
              ]{article}
              
\usepackage[paperheight=297mm,
			paperwidth=210mm,
			tmargin=28mm, 
			headheight=0mm, 
			headsep=15mm,
			textheight=240mm,
			footskip=7mm,
			textwidth=159.2mm,
			bindingoffset=10mm,
			twoside
			]{geometry} 
% -------------------------------------------------------   


% Configuración de párrafos
% -------------------------------------------------------         
\setlength{\parindent}{0em} % Tabulación de un párrafo nuevo
\setlength{\parskip}{1em}   % Distancia en blanco entre un párrafo y otro.
\setlength{\columnsep}{1em} % Margen izquierdo de columnas
\renewcommand{\baselinestretch}{1.5} % Interlineado. Espacio entre dos renglones.
% ------------------------------------------------------- 


% Paquetes de idioma
\usepackage[utf8]{inputenc}
\usepackage[english,spanish,es-tabla]{babel} % English para el abstract
% ------------------------------------------------------- 

% Paquetes para imprimir código fuente
%\usepackage[cache=false]{minted}

\usepackage{subcaption}

\usepackage[table,xcdraw]{xcolor} % Color de tabla

\usepackage{lscape} % Para poder poner en horizontal una página

% Para símbolos de comprobación y x.
\usepackage{pifont}% http://ctan.org/pkg/pifont
\newcommand{\cmark}{ \color{green} \ding{51}}%
\newcommand{\xmark}{ \color{red}\ding{55}}%

\usepackage{eurosym} % Para símbolo €
\usepackage{amsmath} % Para símbolo € en fórmulas

% Paquetes para tree. Gráfico jerárquico
\usepackage{tikz}
\usetikzlibrary{trees}

% -------------------------------------------------------  

% Paquetes generales
\usepackage{enumitem} % Para configurar enumeraciones
%\usepackage[usenames,dvipsnames]{xcolor}    % Para colores en el texto

% Lorem ipsum
\usepackage[pangram]{blindtext}

\usepackage{scrextend} % Para entorno{addmargin} (dar margen a un párrafo...)

% Lorem ipsum
\usepackage{lipsum}

\usepackage{fancyhdr}

%\usepackage[toc,page]{appendix}

\usepackage[spanish]{babel}
\usepackage{appendix}

\renewcommand{\appendixname}{Anexos}
\renewcommand{\appendixtocname}{Anexos}
\renewcommand{\appendixpagename}{Anexos}


\usepackage{float} % para controlar la situación de los entornos flotantes

\usepackage{graphicx} % figuras

% Definimos un nuevo color
\definecolor{webred}{rgb}{0.5, 0, 0}   % less intense red

\usepackage[bookmarks=true,
			bookmarksnumbered=false, % true means bookmarks in 
									 % left window are numbered                         
			bookmarksopen=false,     % true means only level 1
									 % are displayed.
			colorlinks=true,
			linkcolor=webred,
			%linkcolor=OliveGreen,
			citecolor=webred,
			urlcolor=cyan]{hyperref} % Para hiperenlaces

% Definimos una nueva orden para pintar una línea cuando escribamos \HRule
\newcommand{\HRule}{\rule{\linewidth}{1mm}}




% ###################   FIN CONFIGURACIÓN ################### 

% ###################	INICIO DEL DOCUMENTO ###################

\begin{document}

\pagestyle{empty}

	\begin{figure}[h]
		\centering
		\includegraphics[width=0.7\textwidth]{images/logoUGR.jpg}
		\label{imagen2}
	\end{figure}

	\begin{center}
		\Large{TRABAJO FIN DE MÁSTER} \\
		\large{INGENIERÍA INFORMÁTICA}
	\end{center}
	
	\vspace{0.5cm}
	
	\hrule
	
	\begin{center}
		\huge{\textbf{Breaking Transposition Cipher with Genetic Algorithm}} \\
			\end{center}
	
	\hrule
	
	\vspace{0.5cm}
	
	\begin{center}
		\textbf{Autor} \\
		Abdullah Taher Saadoon ALMusawi  (alumno)\\
		\textbf{Directores}\\
		Juan Julián Merelo Guervós (tutor)
	\end{center}
	
	\vspace{0.6cm}
	
	\begin{figure}[h]
		\centering
		\includegraphics[width=0.3\textwidth]{images/etsiit_logo.png}\\[0.1cm]
	\end{figure}
	
	\begin{center}
	\textsc{Escuela Técnica Superior de Ingenierías Informática y de Telecomunicación}\\
	\textsc{---}\\
	Granada, septiembre de 2020
	\end{center}
	
	\newpage

\thispagestyle{empty}

\section*{Acknowledgments}



\newpage

\thispagestyle{empty}

\vspace*{\fill} % Para alinear el contenido de la paǵina verticalmente.

\begin{center}
	\textbf{Breaking Transposition Cipher with Genetic Algorithm}\\
	
	\vspace{0.3cm}
	
	Abdullah Taher Saadoon ALMusawi (alumno)	\\
	
	\footnotesize{Estudiante de Informática en Escuela Técnica Superior Informática y Telecomunicación, Universidad de Granada, 18071, Granada, España} \\
	
\end{center}

\selectlanguage{english} 

\begin{abstract}

\noindent \textbf{Keywords}: Cryptanalysis, Genetic Algorithm, Transposition Cipher, substitution cipher, decryption ,optimization search. 

\vspace{0.2cm}

\noindent 1

\noindent 2

\end{abstract}

\vspace*{\fill} % Para alinear el contenido de la paǵina verticalmente.

\selectlanguage{spanish}

\newpage

\thispagestyle{empty}


\noindent\rule[-1ex]{\textwidth}{2pt}\\[4.5ex]

Yo, \textbf{Abdullah Taher Saadoon AL-Musawi}, alumno de la titulación máster de ingeniería informática de la \textbf{Escuela Técnica Superior de Ingenierías Informática y de Telecomunicación de la Universidad de Granada}, con NIE Y6178000-F, autorizo la ubicación de la siguiente copia de mi Trabajo Fin de Máster en la biblioteca del centro para que pueda ser consultada por las personas que lo deseen.

\vspace{0.5cm}

\noindent Fdo: Abdullah Taher Saadoon AL-Musawi

\begin{figure}[h]
	\centering
	
\end{figure}

\begin{flushright}
Granada a 10 de septimpre de 2020.
\end{flushright}


\newpage

\thispagestyle{empty}

\noindent\rule[-1ex]{\textwidth}{2pt}\\[4.5ex]

D. Juan Julián Merelo Guervos (tutor), Profesor del Departamento de Arquitectura y Tecnología de Computadores de la Universidad de Granada.

\textbf{Informa:}




\textbf{El director:}

\begin{figure}[h]
	\centering
	
\end{figure}





\newpage

\tableofcontents

\newpage



\newpage



\newpage



\newpage



\newpage



\newpage



\newpage



\newpage



\newpage



\end{document}



% ###################	FIN DEL DOCUMENTO ###################